\section{Limitations and Future Work}\label{sec:limitations}

The assumption that inputs and models are benign-but-confusable is foundational to our threat model: covert channels from an actively malicious base model are out of scope, as is malice originating from fully trusted principal sources.

Branching requires a confining runtime. An execution gateway that cannot prevent data generated by a child branch from entering the parent context prior to an explicit merge cannot enforce isolation guarantees, regardless of any labels attached afterward. Duplicating context also incurs token usage, latency, and state-management overheads evaluated in Section~\ref{sec:evaluation}.

Taint prevention enforces parent-local context isolation rather than transactional state rollback. An isolated child branch may execute authorized side effects or external network egresses; discarding the branch frees only its local context and proposed return value. While the shared event log records these external actions for policy evaluation, it cannot roll back state changes in external systems.

Safe branch exit is only as strong as the registered derivation function. Syntactic schema validation cannot prove the absence of semantic influence, and quarantine-exit attestations must be provided by the trusted computing base (TCB). While branching restricts the propagation scope of untrusted taint, it does not eliminate the classical sanitizer-correctness problem.

Contract labels represent declared provenance rather than inferred semantic control flow. APPA does not detect paraphrase, implicit control dependencies, or cross-session data laundering as NeuroTaint does~\cite{caiGhostAgentRedefining2026}, nor does it persist signed derivation lineage across memory as MemLineage does~\cite{ouyangMemLineageLineageGuidedEnforcement2026}. While such mechanisms can inform resolver inputs or audit APPA traces, they remain outside the core algebra.

Approval fatigue represents a primary operational threat to any human-in-the-loop (HITL) system and remains an adoption concern outside the scope of our formal model: the security guarantees of the algebra remain bounded by the correctness of the authorities and contracts registered into it.

\fixme{Dispatch atomicity concerns the crash gap between invoking an external tool and appending its events to the log. Option (a) uses a durable outbox where tool invocation and event entries commit atomically before execution. This storage-layer mechanism is transparent to contract authors and maintains a consistent log. Option (b) splits events by polarity (\texttt{attempted} and \texttt{done}, where restrictive policies gate on attempted events and permissive policies gate on done events). While this fails closed under failure, it reintroduces a pre/post execution axis into every contract and leaves permanent conservative entries in the log.}

\fixme{Future extension topics include outcome-sensitive completion events and branch-scoped \texttt{prior(k)} semantics to prevent a quarantined child branch from generating parent-gating events through benign tool dispatches.}

\paragraph{Epoch-wide authorization scope.}
Ruling-carried persistent authority widenings could in principle be evicted by a subsequent narrowing. This extension is algebraically coherent—encompassing permissive actions, non-commutative folds, the collapse property (Proposition~\ref{prop:collapse}), eviction soundness (\S\ref{sec:enforcement}), and a request-channel design (contract annotations versus alternate execution plans or harness configuration). However, for ongoing interactions with an external recipient, branching provides a superior construction. A child branch confines recipient-specific exchanges and restrictive observations, call-scoped rulings admit only reviewed transmissions, and the child context subsequently terminates. Because \texttt{merge}-by-intersection prevents unauthorized context propagation back to the parent, APPA maintains strict monotone narrowing and focuses implementation complexity on branch isolation and verified exits.

In practice, deployments may attempt to amortize human review using automated approval authorities. However, such an authority must compare each rendered invocation against the specific call it reviewed; failing to perform this comparison degrades security compared to explicit epoch-wide authorization.

Future work includes mandate scoping by information type, pass-by-reference labeling where byte-identical unretyped arguments retain their original labels (a consequence of arguments being \texttt{LabeledValue} instances), and partial-context handoff, in which a child branch is instantiated with a sanitized subset of the parent's context while retaining access to full shared event history.

Policy dimensions must undergo engineering review prior to deployment: while algebraic lawfulness can be property-tested, security semantics cannot. Candidate dimensions include spending budgets as a canonical monoid (ordered, per-run, and accumulating monotonically—currently evaluated via log aggregation in a resolver-implemented authority), independent attestations, and value-specific purpose limitations.

